%%%%%%%%%%%%%%%%%%%%%%%%%%%%%%%%%%%%%%%%%%%%%%%%%%%%%%%%%%%%%%%%%%%%%%%%%%%
% Englische Fassung des Uebungsblatts zu Lehreinheit 4 (Teil A und B).
% Vorlage: handouts/LE04.tex. Nur in handouts.tex mit \def\wuzsprache{en}.
%%%%%%%%%%%%%%%%%%%%%%%%%%%%%%%%%%%%%%%%%%%%%%%%%%%%%%%%%%%%%%%%%%%%%%%%%%%
\begin{uebungsblatt}{Source Criticism and SIFT}

\aufgabe{SIFT grid for three sources}[in class, small groups, 30 min.; finish as homework]
The instructor provides three sources: (1)~a peer-reviewed study, (2)~a
newspaper article, (3)~a Wikipedia entry. First enter the exact details,
then your findings for each step. In step~T, trace at least one number back
to the original study. On the type of source: \emph{primary} = own data
(study), \emph{secondary} = summarises studies (review, meta-analysis),
\emph{tertiary} = reports on them (press, encyclopedia), \emph{grey} =
published without a publisher or peer review (preprint, report, thesis).

\begin{center}
\small
\renewcommand{\arraystretch}{1.3}
\begin{tabular}{|>{\raggedright\arraybackslash}p{0.15\textwidth}|p{0.24\textwidth}|p{0.24\textwidth}|p{0.24\textwidth}|}
\hline
 & \textbf{(1) Study} & \textbf{(2) Newspaper article} & \textbf{(3) Wikipedia} \\
\hline
Details (author, title, date, link) & \rule{0pt}{14mm} & & \\ \hline
\textbf{S}top: What do I want to know? First impression? & \rule{0pt}{16mm} & & \\ \hline
\textbf{I}nvestigate: Who is the source? What do others say? & \rule{0pt}{18mm} & & \\ \hline
\textbf{F}ind: Better or independent coverage? & \rule{0pt}{16mm} & & \\ \hline
\textbf{T}race: Where does the central number/claim come from? & \rule{0pt}{18mm} & & \\ \hline
Type of source (primary, secondary, tertiary, grey) & \rule{0pt}{9mm} & & \\ \hline
Weight (high/medium/low) and reason & \rule{0pt}{16mm} & & \\ \hline
\end{tabular}
\end{center}
\end{uebungsblatt}

\begin{uebungsblattb}{2}{Beautiful Fakes and Source Criticism}

\aufgabe{Anatomy of a beautiful fake}[in class after the reveal, individual work, 10 min.]
The following news item is \textbf{fictitious} and was constructed for
this exercise, following the pattern of the \enquote{beautiful fake} from
the moment of wonder. Mark every feature that \emph{creates} credibility,
and note for each whether it \emph{proves} anything.

\begin{center}
\fbox{\parbox{0.9\textwidth}{\small
\textbf{Juggling makes you smart: ten minutes in the morning boost memory
by 23.7 percent}\\[1mm]
People who start the day with juggling balls learn faster. That is the
finding of a new study that experts have been discussing for weeks.
\enquote{The effects are impressive}, says Prof.\ Dr.\ H.\ Lindner,
learning researcher and author of several self-help books. Participants who
juggled before learning vocabulary retained on average 23.7\,\% more words.
More and more schools are already considering juggling breaks. This fits
what educators have long suspected: movement and thinking belong
together.\\[1mm]
{\footnotesize\itshape (Fictitious teaching material, Unit~4. Person, study and figures are invented.)}}}
\end{center}

\begin{center}
\renewcommand{\arraystretch}{1.9}
\begin{tabular}{|p{0.42\textwidth}|p{0.48\textwidth}|}
\hline
\textbf{Feature} & \textbf{Does it prove anything? What is missing?} \\
\hline
 & \\ \hline
 & \\ \hline
 & \\ \hline
 & \\ \hline
\end{tabular}
\end{center}

\aufgabe{Source criticism in five steps}[homework]
During the wrap-up, enter the five steps that the group works out on the
board in the left column (for comparison: handout). Then apply the
checklist to two entries of your annotated bibliography from Unit~3. Tick
what you have checked and note whether your judgement has changed.

\begin{center}
\renewcommand{\arraystretch}{1.6}
\begin{tabular}{|p{0.5\textwidth}|c|c|}
\hline
\textbf{Step} & \textbf{Source A} & \textbf{Source B} \\
\hline
1 \rule{0pt}{7mm} & $\square$ & $\square$ \\ \hline
2 \rule{0pt}{7mm} & $\square$ & $\square$ \\ \hline
3 \rule{0pt}{7mm} & $\square$ & $\square$ \\ \hline
4 \rule{0pt}{7mm} & $\square$ & $\square$ \\ \hline
5 \rule{0pt}{7mm} & $\square$ & $\square$ \\ \hline
\end{tabular}
\end{center}
Has your judgement changed? Why?
\schreiblinien{3}

\aufgabe{Discussion: 500 citations}[homework, max.\ 1 page]
\enquote{A study has 500 citations. What does that say about its quality?}
Write a structured counter-argument against the assumption that many
citations mean high quality. Structure: (a)~the thesis in its strongest
form, (b)~at least three counter-arguments, each with evidence or an
example, (c)~what citations \emph{do} tell us, (d)~your conclusion: how
would you assess quality instead?

\aufgabe{Diary of wonder}[homework]
Describe a pattern you noticed this week (chance, statistics, news,
everyday life). What did you expect, what did you see, and how could you
test whether the pattern is real?

\begin{lehrhinweis}[Expected answers]
Task~1: see the teaching note in class; good answers name an
\emph{external} source in step~I and, in step~T, the place in the study
where the number comes from. Task~2: credibility features are the precise
number (23.7\,\%), the professor's title, the quotation, \enquote{a new
study} without details, \enquote{experts are discussing}, \enquote{more
and more schools}, the link to what is familiar (\enquote{what educators
have long suspected}). None of them proves anything: study, sample,
control group, measurement procedure and the expert's field are missing
(self-help author: conflict of interest). Task~3: the left column contains
the five steps of the wrap-up (stop; origin and intention; type and method;
trace to the original; cross-check), as in the handout. Task~4: expected
counter-arguments: citations measure attention, not correctness; criticism
and refutation also count as citations; Matthew effect
\parencite{merton1968matthew}; field and age effects (older papers and
large fields collect more citations); self-citations; citations are often
copied unread. It should be conceded that a much-cited paper is present in
the scholarly debate and is usually a good starting point. Stronger answers
propose ways of checking: read the method, look for replications and
retractions, check in what sense the paper is cited
\parencite{dora2012,hicks2015leiden}.
\end{lehrhinweis}
\end{uebungsblattb}
